% =============================================================================
% 5.1 Discussion of SQ1 — The VM Failure Threshold
% =============================================================================

\section{Discussion of SQ1 --- The VM Failure Threshold}
\label{sec:discuss-sq1}

\textbf{SQ1 asks:} At what concurrency level do Docker containers on a single VM degrade in reliability and performance, and how do success rate, execution time variance, and resource utilisation degrade as concurrent load increases toward that threshold?

\paragraph{The reliability threshold was reached at L4.}
Contrary to the pre-experiment hypothesis (which predicted reliability collapse at L5--L7), the VM's success rate dropped sharply at five concurrent jobs: from 100\,\% at L1--L3 to 66.7\,\% at L4.
The failure mechanism is memory exhaustion: with five containers each allocating 400\,MB simultaneously, aggregate memory demand exceeds the available headroom in the 4\,GB VM (which also runs the Docker daemon, PostgreSQL, and Dagster daemon processes consuming $\approx$1.5--2\,GB).
Individual run records confirm that failing jobs are OOM-killed midway through execution --- immediately after the 400\,MB NumPy allocation in the \texttt{memory\_pressure} op.

The Linux OOM killer selects victims based on process memory scores; in a pool of identical containers, the selection is approximately random across the $n$ concurrent runs.
This explains the observed pattern: at L4, one of five containers is typically killed, yielding a 66.7\,\% success rate; at L6, all ten are typically killed, yielding 0\,\%.

\paragraph{The reliability crossover falls between L3 and L4.}
The formal reliability crossover (success rate falling below 95\,\%) occurs between L3 (100\,\%) and L4 (66.7\,\%), at the point where aggregate memory allocation first exceeds the kernel's reclaimable budget on this 4\,GB VM.
This is the most practically significant finding of SQ1: the VM degradation is not a \emph{gradual} performance curve but a \emph{discrete transition} from reliable to unreliable operation at five concurrent jobs.

\paragraph{Execution time of surviving jobs remains stable until L4.}
Successful jobs at L1--L3 complete in approximately 66--69\,s, but then jump dramatically to 168.6\,s at L4 and 464.1\,s at L5.
This is a selection effect combined with resource contention: only containers that successfully acquire their memory allocation complete, and those at higher concurrency experience severe scheduling delays due to extreme memory pressure.
OOM-killed containers are accounted in the failure rate, not in the execution time metric.

\paragraph{Comparison to the hypothesis.}
The hypothesis expected gradual CPU-driven performance degradation as the primary failure mode, with memory-driven OOM kills as a secondary effect at very high concurrency.
The experimental data partially confirmed this: OOM-driven failures dominate from L4 onward, but the sharp three-job transition to unreliability is earlier than predicted (L4 vs. L5--L7 hypothesis).
This reflects the workload design: the 400\,MB allocation per job creates a hard per-container memory footprint that the 4\,GB VM cannot sustain at 5 or more concurrent jobs while also running its infrastructure stack.
This confirms the prediction of \citet{arora2023} that contention under shared execution models produces non-linear failure, though the failure mode is memory-driven rather than CPU-driven.
